\section{Kết luận}
\label{sec:conclusion}

Báo cáo đã phân tích chuỗi lưu lượng dòng chảy sông Colorado giai đoạn 2010--2023
và trả lời ba nhóm câu hỏi đặt ra. Về cấu trúc chuỗi, dữ liệu có \DatasetN{} ngày,
không thiếu, không trùng và không có giá trị bất thường về miền xác định. Chuỗi có
tính dai dẳng ngắn hạn rất mạnh, mùa vụ theo tháng yếu và mức nền thay đổi theo
từng giai đoạn. Xu hướng STL không tuyến tính: mức nền tăng từ đầu giai đoạn, đạt
đỉnh ở giữa và giảm trong những năm gần đây. Đây là đặc điểm quan trọng nhất cần
được tính đến khi chọn mô hình.

Về cực trị, ngưỡng Q95 là \QHigh{} m\textsuperscript{3}/s, tương ứng z tĩnh
\ZThreshold{}. Có \NDaysHigh{} ngày vượt ngưỡng, chiếm \ShareHigh{}\%, và
\NEvents{} đợt. Các đợt tập trung vào 2017--2018 và 2020; đợt dài nhất kéo dài
\LongestEvent{} ngày. Phân bố theo tháng không tạo thành một mùa vụ cố định, và
ngưỡng toàn mẫu chỉ phù hợp cho mục tiêu mô tả, không phù hợp cho vận hành.

Về dự báo, LightGBM direct 7 đạt mean RMSE \LgbmMeanRmse{} và NSE \LgbmMeanNse{}
trên Test, gần tương đương Naive (\NaiveMeanRmse{} và \NaiveMeanNse{}) và tốt hơn
nhiều so với Seasonal Naive 365 (\SeasonalMeanRmse{}). Khoảng cách giữa LightGBM và
Naive chỉ khoảng 0,3\%, không đủ để khẳng định mô hình học máy vượt trội. Sai số
tăng theo horizon; D+1 có RMSE khoảng 2,79 và NSE 0,970, trong khi D+7 có RMSE
khoảng 7,10 và NSE 0,802. Test không có ngày vượt Q95, nên báo cáo không thể kết
luận về khả năng dự báo cực trị.

Ba bài học chính rút ra là: (i) baseline đơn giản phải được đánh giá nghiêm túc
trước khi thêm mô hình phức tạp; (ii) mùa vụ yếu và mức nền thay đổi có thể khiến
các mô hình dựa trên “cùng ngày năm trước” thất bại; (iii) chất lượng dữ liệu và
khả năng tái lập quan trọng không kém độ phức tạp của mô hình. Website kết quả,
mã nguồn và manifest được công bố để người đọc có thể kiểm tra lại từng con số.

Hướng phát triển tiếp theo là bổ sung biến ngoại sinh, đánh giá xác suất, dùng
ngưỡng cực trị trượt và mở rộng sang nhiều trạm. Trước khi làm những việc đó, cần
giữ nguyên nguyên tắc: mọi cải thiện phải được chứng minh trên dữ liệu ngoài mẫu,
với cùng split và cùng baseline. Trong phạm vi dữ liệu hiện có, kết luận trung thực
nhất là LightGBM chưa vượt được Naive, và báo cáo trình bày kết quả đó thay vì che
giấu bằng một phép so sánh có lợi.
